\appendix

\section{Algebras over groupoids in algebraic spaces}\label{appendix}

In this appendix, we prove a generalization of Proposition~\ref{quotient correspond}. We freely use the terminology and notation from~\cite{stacks}, and categories fibered in groupoids are defined over the big fppf site $(\Sch/S)_{\fppf}$ over a fixed base scheme $S$. First, we recall some definitions from~\cite{stacks}, where the main references are Chapters {\it Groupoids in Algebraic Spaces, Algebraic Stacks}, and {\it Sheaves on Algebraic Stacks}.

Let $p\colon \cX\to (\Sch/S)_{\fppf}$ be a category fibered in groupoids. Recall that we have the induced fppf topology on $\cX$, where a family of morphisms $\{x_i\to x\}_{i}$ in $\cX$ is a {\it covering} of $x\in \cX$ if the family $\{ p(x_i)\to p(x)\}_i$ is an fppf cover of the scheme $p(x)$. We write $\cX_{\fppf}$ for this fppf site. Then we have the {\it structure sheaf} $\cO_{\cX}\colon \cX_{\fppf}^{\op}\to (\Rings)$ of $\cX_{\fppf}$ defined by
\begin{gather*}
\cO_{\cX}(x)\defeq\Gamma\big(p(x),\cO_{p(x)}\big),
\end{gather*}
and thus we have the associated ringed site $\big(\cX_{\fppf},\cO_{\cX}\big)$. An $\cO_{\cX}$-module $\cF\colon\cX_{\fppf}^{\op}\to (\Sets)$ is said to be {\it quasi-coherent}, if for any $x\in \cX$ there is a covering $\{x_i\to x\}$ of $x$ such that for each~$x_i$ there is an exact sequence of $\cO_{\cX/x_i}$-modules
\begin{gather*}
\bigoplus_{I}\cO_{\cX/x_i}\to \bigoplus_{J}\cO_{\cX/x_i}\to\cF|_{x_i}\to 0,
\end{gather*}
 where $\big(\cX_{\fppf}/x_i,\cO_{\cX/x_i}\big)$ is the localization of the ringed site $\big(\cX_{\fppf},\cO_{\cX}\big)$ at $x_i\in\cX$ and $\cF|_{x_i}$ is the restriction of $\cF$ to $\big(\cX_{\fppf}/x_i,\cO_{\cX/x_i}\big)$. We denote by $\Qcoh \cX$ the category of quasi-coherent $\cO_{\cX}$-modules. If $\cX$ is an algebraic stack, $\Qcoh \cX$ is equivalent to the category of quasi-coherent sheaves on the lisse-\'etale site of $\cX$ defined in~\cite[Definition 9.1.14]{olsson}. In particular, if $\cX$ is an algebraic space, $\Qcoh \cX$ is equivalent to the category of quasi-coherent sheaves on the small \'etale site of $\cX$ defined in~\cite[Definition 7.1.5]{olsson}, and if $\cX$ is a scheme, it is equivalent to the category of usual quasi-coherent sheaves on the small Zariski site of $\cX$. If $X$ is an algebraic space, we write $X_{\et}$ (resp.\ $X_{\fppf}$) for the small \'etale site of $X$ (resp.\ the big fppf site $(\Sch/X)_{\fppf}$ of $X$), and sheaves on $X$ means sheaves on $X_{\et}$. Similarly, sheaves on a scheme $X$ means sheaves on the small Zariski site of $X$.


 \begin{dfn} An {\it ${\cX}$-algebra} is a sheaf of (not necessarily commutative) rings $\cA$ on $\cX_{\fppf}$ together with a morphism of sheaves of rings
 \begin{gather*}
 \uprho\colon\ \cO_{\cX}\to \cA
 \end{gather*}
 such that the $\cO_{\cX}$-module $\cA$ is quasi-coherent and the image of $\uprho$ is in the center of $\cA$, i.e., for any $x\in \cX$ the image of $\uprho(x)\colon \cO_{\cX}(x)\to \cA(x)$ is contained in the center of the ring $\cA(x)$.
\end{dfn}

 If $(\cA,\uprho)$ is an $\cX$-algebra, we have the restriction by $\uprho$
 \begin{gather*}
 (-)_{\uprho}\colon\quad \Mod\cA\to \Mod\cO_{\cX}, \qquad \cM\mapsto \cM_{\uprho},
 \end{gather*}
 where $ \Mod\cA$ denotes the category of right $\cA$-modules.
 \begin{dfn} A right $\cA$-module $\cM$ is said to be {\it quasi-coherent} if the $\cO_{\cX}$-module $\cM_{\uprho}$ is quasi-coherent. We denote by
 \begin{gather*}
 \Qcoh\cA
 \end{gather*}
 the full subcategory of $\Mod\cA$ consisting of quasi-coherent right $\cA$-modules.
 \end{dfn}

 \begin{rem}
 If $\cX$ is an algebraic space, we tacitly consider $\cX$-algebras and quasi-coherent modules over $\cX$-algebras as sheaves on the small \'etale site $\cX_{\et}$ instead of big fppf sheaves as above.
 \end{rem}



Let $\cA$ be an $\cX$-algebra, and denote by $\widetilde{\cX}$ the stackification of $\cX$. By~\cite[Lemma~12.1; 06WQ]{stacks}, we have an equivalence of topoi
\begin{equation}\label{equiv topoi}
\Sh(\cX_{\fppf})\simto \Sh\big(\widetilde{\cX}_{\fppf}\big),
\end{equation}
where $\Sh(-)$ denotes the category of sheaves on a site $(-)$. For an object $\cF\in \Sh(\cX_{\fppf})$, we denote by $\widetilde{\cF}\in\Sh\big(\widetilde{\cX}_{\fppf}\big)$ the image of $\cF$ by the equivalence \eqref{equiv topoi}.
By this equivalence the ring object $\cA\in\Sh(\cX_{\fppf})$ defines a ring object $\widetilde{\cA}\in\Sh\big(\widetilde{\cX}_{\fppf}\big)$ and a morphism
\begin{gather*}
\widetilde{\uprho}\colon\ \cO_{\widetilde{\cX}}\to \widetilde{\cA}.
\end{gather*}
By~\cite[Lemma~12.2; 06WR]{stacks}, the sheaf of rings $\widetilde{\cA}$ is quasi-coherent $\cO_{\widetilde{\cX}}$-module, and so $\widetilde{\cA}$ is an~$\widetilde{\cX}$-algebra.




\begin{lem}\label{stackification equiv}
The equivalence of topoi $\Sh(\cX_{\fppf})\simto \Sh\big(\widetilde{\cX}_{\fppf}\big)$ induces the equivalence
\begin{gather*}
\Mod\cA\simto \Mod\widetilde{\cA}
\end{gather*}
of right modules, and it restricts to the equivalence of quasi-coherent modules
\begin{gather*}
\Qcoh\cA\simto \Qcoh\widetilde{\cA}.
\end{gather*}
\begin{proof}
The first assertion is obvious, and for the second statement it is enough to show that a right $\cA$-module $\cM$ is quasi-coherent if and only if the right $\widetilde{\cA}$-module $\widetilde{\cM}$ is quasi-coherent. But this follows from the following commutative diagram
\[
\begin{tikzcd}
\Mod\cA\arrow[rr,"\widetilde{(-)}"]\arrow[d,"(-)_{\uprho}"']&&\Mod\widetilde{\cA}\arrow[d," (-)_{\widetilde{\uprho}}"]\\
\Mod\cO_{\cX}\arrow[rr, "\widetilde{(-)}"]&&\Mod\cO_{\widetilde{\cX}}
\end{tikzcd}
\]
and~\cite[Lemma~12.2; 06WR]{stacks}.
\end{proof}
\end{lem}


Let $\cG=(U,R,s,t, c)$ be a groupoid in algebraic spaces over $S$ (see~\cite[Definition 11.1; 043W]{stacks} for the notation).


\begin{dfn}
 An {\it algebra over $\cG$}, or {\it $\cG$-algebra}, is a $U$-algebra $\cA\colon U_{\et}^{\op}\to (\Rings)$ together with an isomorphism
\begin{gather*}
\uptheta^{\cA}\colon\ s^*\cA\simto t^*\cA
\end{gather*}
of $R$-algebras such that the equations
\begin{equation}\label{groupoid alg}
p_1^*\uptheta^{\cA}\circ p_2^*\uptheta^{\cA}=c^*\uptheta^{\cA}
\qquad\text{and}\qquad
e^*\uptheta^{\cA}=\id_{\cA}
\end{equation}
of morphisms of sheaves of rings hold, where $p_i\colon R\times_{s,U,t}R \to R$ is the $i$-th projection and $e\colon U\to R$ is the identity.
\end{dfn}

\begin{rem}\label{str equiv}
Note that the structure sheaf $\cO_U$ together with the canonical isomorphism
\begin{gather*}
\uptheta^{\can}\colon\ s^*\cO_U\simto t^*\cO_U
\end{gather*}
defines a $\cG$-algebra $(\cO_U, \uptheta^{\can})$. If $\big(\cA,\uptheta^{\cA}\big)$ is a $\cG$-algebra, since $\uptheta^{\cA}$ is $\cO_R$-linear, the diagram
\[
\begin{tikzcd}
s^*\cO_U\arrow[rr,"\sim"]\arrow[d,"s^*\uprho"']\arrow[rrrr, bend left=20,swap, "\uptheta^{\can}"']&&\cO_R\arrow[lld," "]\arrow[rrd," "]\arrow[rr,"\sim"]&&t^*\cO_U\arrow[d,"t^*\uprho"]\\
s^*\cA\arrow[rrrr, "\uptheta^{\cA}"]&&&&t^*\cA
\end{tikzcd}
\]
is commutative, where the isomorphisms on the top level are natural isomorphisms.
\end{rem}


\begin{dfn} Notation is the same as above.
\begin{itemize}\itemsep=0pt
\item[1.] A {\it quasi-coherent module} over a $\cG$-algebra $\big(\cA,\uptheta^{\cA}\big)$ is a quasi-coherent right $\cA$-module $\cM\colon U^{\op}_{\et}\to (\Sets)$ together with an isomorphism
\begin{gather*}
\uptheta^{\cM}\colon\ s^*\cM\simto t^*\cM
\end{gather*}
of quasi-coherent right $s^*\cA$-modules such that the similar equations as \eqref{groupoid alg} hold, where the right $t^*\cA$-module $t^*\cM$ is considered as a right $s^*\cA$-module via the isomorphism $\uptheta^{\cA}\colon s^*\cA\simto t^*\cA$ of $R$-algebras.
\item[2.] Let $\big(\cM,\uptheta^{\cM}\big)$ and $\big(\cN,\uptheta^{\cN}\big)$ be quasi-coherent modules over a $\cG$-algebra $\big(\cA,\uptheta^{\cA}\big)$.
A {\it morphism} from $\big(\cM,\uptheta^{\cM}\big)$ to $\big(\cN,\uptheta^{\cN}\big)$ is a morphism
\begin{gather*}
\varphi\colon\ \cM\to \cN
\end{gather*}
of right $\cA$-modules such that the diagram
\[
\begin{tikzcd}
s^*\cM\arrow[rr,"s^*\varphi"]\arrow[d,"\uptheta^{\cM}"']&&s^*\cN\arrow[d," \uptheta^{\cN}"]\\
t^*\cM\arrow[rr, "t^*\varphi"]&&t^*\cN
\end{tikzcd}
\]
is commutative.
\end{itemize}
\end{dfn}

We denote by
\begin{gather*}
\Qcoh\big(\cA,\uptheta^{\cA}\big)
\end{gather*}
the category of quasi-coherent modules over a $\cG$-algebra $\big(\cA,\uptheta^{\cA}\big)$.

\begin{rem}
We have a natural identification
\begin{gather*}
\Qcoh\big(\cO_U,\uptheta^{\can}\big)=\Qcoh(U,R,s,t,c),
\end{gather*}
where $\Qcoh(U,R,s,t,c)$ denotes the category of quasi-coherent modules on $(U,R,s,t,c)$ in the sense of~\cite[Definition 12.1; 0441]{stacks}.
\end{rem}


Let $\cG=(U,R,s,t, c)$ be a groupoid in algebraic spaces over $S$. Following~\cite[044O]{stacks} we denote by $[U/_pR]$ the category fibered in groupoids associated to $\cG$, which is defined as follows: An object of $[U/_pR]$ is a pair $(T,u)$ of $S$-scheme $T$ and a morphism $u\colon T\to U$ of algebraic spaces over $S$, and a morphism $(f,\varphi)\colon(T,u)\to (T',u')$ is defined to be a pair of a morphism $f\colon T\to T'$ of $S$-schemes and a morphism $\varphi\colon T\to R$ of algebraic spaces over $S$ such that $s\circ \varphi=u$ and $t\circ \varphi=u'\circ f$. Then the functor
\begin{gather*}
[U/_pR]\to (\Sch/S), \qquad (T,u)\mapsto T
\end{gather*}
is a category fibered in groupoids.
We denote by $[U/R]$ the stackification of $[U/_pR]$, which is called the {\it quotient stack} of $\cG$. We have the natural morphism
\begin{equation}\label{atlas}
\alpha\colon\ (\Sch/U)\to [U/_pR]
\end{equation}
defined by $\alpha(u\colon T\to U)\defeq(T,u)$ and $\alpha(f\colon T\to T')\defeq(f, e\circ u)$, and it defines the functor
\begin{gather*}
\alpha^{-1}\colon\ \Sh\big([U/_pR]_{\fppf}\big)\to \Sh(U_{\fppf})
\end{gather*}
defined by $\alpha^{-1}\cF(u\colon T\to U)\defeq \cF(\alpha(u\colon T\to U))=\cF(T,u)$.

For a $\cG$-algebra $\big(\cA,\uptheta^{\cA}\big)$, we define a sheaf of rings
\begin{gather*}
[\cA/_p R]\colon\ [U/_pR]_{\fppf}^{\op}\to (\Rings)
\end{gather*}
on $[U/_pR]_{\fppf}$ as follows:
For an object $(T,u)\in [U/_pR]$ we set
\begin{gather*}
[\cA/_p R](T,u)\defeq\Gamma(T,u^*\cA),
\end{gather*}
%(where, by abuse of notation, we write $\cA$ for the big fppf sheaf associated to the small \'etale sheaf $\cA$,)
and for a morphism $(f,\varphi)\colon (T,u)\to (T',u')$ we define a morphism of rings
\begin{gather*}
[\cA/_p R](f,\varphi)\colon\ [\cA/_p R](T',u')\to [\cA/_p R](T,u)
\end{gather*}
by the following commutative diagram:
\[
\begin{tikzcd}
\Gamma(T',u'^*\cA)\arrow[rrr,"{[\cA/_p R](f,\varphi)}"]\arrow[d,"f^*"']&&&\Gamma(T,u^*\cA)\\
\Gamma(T,f^*u'^*\cA)\arrow[r, "\sim"]&\Gamma(T,\varphi^*t^*\cA)\arrow[rr,"\varphi^*{\uptheta^{\cA}}"]&&\Gamma(T,\varphi^*s^*\cA)\arrow[u, "\sim"' labl].
\end{tikzcd}
\]

 Similarly, for a quasi-coherent module $\big(\cM,\uptheta^{\cM}\big)$ over a $\cG$-algebra $\big(\cA,\uptheta^{\cA}\big)$, we define the right $[\cA/_pR]$-module
\begin{gather*}
[\cM/_pR]\colon\ [U/_pR]_{\fppf}^{\op}\to (\Sets)
\end{gather*}
by $[\cM/_p R](T,u)\defeq\Gamma(T,u^*\cM)$, and this defines an additive functor
\begin{equation}\label{quot functor}
[(-)/_pR]\colon\ \Qcoh\big(\cA,\uptheta^{\cA}\big)\to \Mod[\cA/_pR].
\end{equation}
By the equality $[\cO_U/_pR]=\cO_{[U/_pR]}$ and Remark~\ref{str equiv}, the morphism $\uprho\colon \cO_U\to \cA$ defines a morphism from $(\cO_U,\uptheta^{\can})$ to $\big(\cA,\uptheta^{\cA}\big)$ in $\Qcoh\big(\cA,\uptheta^{\cA}\big)$, and thus we have the induced morphism
\begin{gather*}
[\uprho/_pR]\colon\ \cO_{[U/_pR]}\to [\cA/_pR]
\end{gather*}
of sheaves of rings.

\begin{lem}\label{qcoh qcoh} Notation is the same as above.
\begin{itemize}\itemsep=0pt
\item[$1.$] If $(\cM,\uptheta^{\can})\in \Qcoh(\cO_U,\uptheta^{\can})$, then $[\cM/_pR]\in \Qcoh[U/_pR]$.
\item[$2.$] The sheaf $[\cA/_pR]$ together with the ring homomorphism $[\uprho/_pR]$ is a $[U/_pR]$-algebra.
\item[$3.$] If $\big(\cM,\uptheta^{\cM}\big)\in \Qcoh\big(\cA,\uptheta^{\cA}\big)$, the right $[\cA/_pR]$-module $[\cM/_pR]$ is quasi-coherent.
\end{itemize}
\begin{proof}
(1) Let $\{a_i\colon A_i\to U\}_{i\in I}$ be an affine \'etale covering of $U$. Since $\cM$ is quasi-coherent $\cO_U$-module, for each $i\in I$ there is an exact sequence
\begin{equation}\label{qcoh exact}
\bigoplus_{J}\cO_{A_i}\to \bigoplus_{K}\cO_{A_i}\to a_i^*\cM\to 0.
\end{equation}
For any object $(T,u)\in [U/_pR]$, set $T_i\defeq A_i\times_{U}T$ and write $t_i\colon T_i\to T$ for the second projection. If we set $u_i\defeq u\circ t_i\colon T_i\to U$ and $\varphi_i\defeq e\circ u_i\colon T_i\to R$, then
\begin{gather*}
\bigl\{(t_i,\varphi_i)\colon (T_i, u_i)\to (T,u)\bigr\}_{i\in I}
\end{gather*}
is a covering of $(T,u)\in [U/_pR]_{\fppf}$. Since the localized ringed site
\begin{gather*}
\big([U/_pR]_{\fppf}/(T_i,u_i),\cO_{[U/_pR]_{\fppf}/(T_i,u_i)}\big)
\end{gather*}
is equivalent to the ringed site $\big(({\Sch/T_i})_{\fppf},\cO_{({\Sch/T_i})_{\fppf}}\big)$ by~\cite[Lemma~9.4; 06W9]{stacks}, the pull-back of the exact sequence \eqref{qcoh exact} by the first projection $T_i\to A_i$ gives rise to an exact sequence
\begin{gather*}
\bigoplus_{J}\cO_{[U/_pR]_{\fppf}/(T_i,u_i)}\to \bigoplus_{K}\cO_{[U/_pR]_{\fppf}/(T_i,u_i)}\to \cM|_{(T_i,u_i)}\to 0.
\end{gather*}

(2) Since the image of $\uprho\colon \cO_U\to \cA$ is contained in the center of $\cA$, the image of \begin{gather*}[\uprho/_pR]\colon \cO_{[U/_pR]}\to [\cA/_pR]\end{gather*} is contained in the center of $[\cA/_pR]$. The $\cO_{[U/_pR]}$-module $[\cA/_pR]$ is quasi-coherent by (1).

(3) This follows from (1).
\end{proof}
\end{lem}

By Lemma~\ref{qcoh qcoh}, the functor \eqref{quot functor} defines an additive functor
\begin{gather*}
[(-)/_pR]\colon\ \Qcoh\big(\cA,\uptheta^{\cA}\big)\to \Qcoh[\cA/_pR].
\end{gather*}
\begin{prop}\label{quot equiv}
The functor $[(-)/_pR]\colon \Qcoh\big(\cA,\uptheta^{\cA}\big)\to \Qcoh[\cA/_pR]$ is an equivalence.
\begin{proof}
This follows from a similar argument as in the proof of~\cite[Proposition~13.1; 06WT]{stacks}. We construct a quasi-inverse
of the functor $[(-)/_pR]$ as follows.

First, for a sheaf $\cF$ on $[U/_pR]_{\fppf}$, we define a sheaf $\cF^R$ on the small \'etale site $U_{\et}$ by
\begin{gather*}
\cF^R\defeq\big(\alpha^{-1}\cF\big)|_{U_{\et}},
\end{gather*}
where $\alpha\colon (\Sch/U)\to [U/_pR]$ is the morphism in \eqref{atlas}. By construction we have $\cO_{[U/_pR]}^R=\cO_{U_{\et}}$. For each \'etale morphism $(u\colon T\to U)\in U_{\et}$, we have
\begin{gather*}
[\cA/_pR]^R(u\colon T\to U)=\Gamma(T,u^*\cA)=\cA(u\colon T\to U).
\end{gather*}
 Thus we have a natural identification $[\cA/_pR]^R=\cA$, and so if $\cM$ is a quasi-coherent right $[\cA/_pR]$-module, the small \'etale sheaf $\cM^R$ is a quasi-coherent $\cA$-module.

Next, we define a category $\{ U/_pR\}$ containing the category $[U/_pR]$ by replacing $S$-sche\-mes~$T$ in the definition of $[U/_pR]$ with algebraic spaces over $S$. If $(f,\varphi)\colon (W,u)\to(W',u')$ is a~morphism in $\{U/_pR\}$, this morphism induces a natural (functorial) commutative diagram
\[
\begin{tikzcd}
W \arrow[rr,"f"]\arrow[rd,"\alpha\circ u"']&&W'\arrow[ld," \alpha\circ u'"]\\
&\left[U/_pR\right].&
\end{tikzcd}
\]
Then, for any quasi-coherent right $[\cA/_pR]$-module $\cM$, we have the composition
\begin{gather*}
(\alpha\circ u')^{-1}\cM\to f_*f^{-1}\big((\alpha\circ u')^{-1}\cM\big)\simto f_*(\alpha\circ u)^{-1}\cM,
\end{gather*}
where the first morphism is the adjunction morphism and the second isomorphism follows from the above commutative diagram. By the adjunction $f^{-1}\dashv f_*$, this composition map induces the morphism
$f^{-1}(\alpha\circ u')^{-1}\cM\to (\alpha\circ u)^{-1}\cM$
and restricting this morphism to $W_{\et}$, we have the comparison map
\begin{gather*}
c_{(f,\varphi)}\colon\ f^*u'^*\cM^R\to u^*\cM^R
\end{gather*}
and by~\cite[Lemma~11.6; 06WK]{stacks} this is an isomorphism.

Finally, we define a functor $F\colon \Qcoh[\cA/_pR]\to \Qcoh\big(\cA,\uptheta^{\cA}\big)$ by
\begin{gather*}
F(\cM)\defeq\big(\cM^R, \uptheta^{\cM^R}\big),
\end{gather*}
where $\uptheta^{\cM^R}\colon s^*\cM^R\simto t^*\cM$ is defined to be the inverse of the comparison map
\begin{gather*}
c_{(\id_R,\id_R)}\colon\ t^*\cM^R\simto s^*\cM^R
\end{gather*}
induced by the morphism $(\id_R,\id_R)\colon (R,s)\to (R,t)$ in $\{U/_pR\}$. By construction, $F$ is a quasi-inverse of the functor $[(-)/_pR]$.
\end{proof}
\end{prop}

Write
\begin{gather*}
[\cA/R]\defeq\widetilde{[\cA/_pR]}
\end{gather*}
for the sheaf of rings on the quotient stack $[U/R]$ corresponding to $[\cA/_pR]$ via the equivalence~\eqref{equiv topoi}. By Lemma~\ref{stackification equiv} and Proposition~\ref{quot equiv}, we have the following
 generalization of~\cite[Proposition~13.1; 06WT]{stacks}.

\begin{prop}\label{main app}
We have an equivalence
\begin{gather*}
\Qcoh\big(\cA,\uptheta^{\cA}\big)\simto \Qcoh[\cA/R].
\end{gather*}
\end{prop}

